%%% Document Formatting
\documentclass[12pt,fleqn]{article}
\usepackage[a4paper,
            bindingoffset=0.2in,
            left=0.75in,
            right=0.75in,
            top=0.8in,
            bottom=0.8in,
            footskip=.25in]{geometry}
\setlength\parindent{10pt} % No indent

%%% Imports
% Mathematics
\usepackage{amsmath} % Math formatting
\numberwithin{equation}{section} % Number equation per section


\newtheorem{theorem}{Theorem}
\newtheorem{definition}{Definition}
\newtheorem{proof}{Proof}
\newtheorem{lemma}{Lemma}
\newtheorem{example}{Example}

\usepackage{amsmath}
\usepackage{amsfonts} % Math fonts
\usepackage{amssymb} % Math symbols
\usepackage{mathtools} % Math etc.
\usepackage{slashed} % Dirac slash notation
\usepackage{cancel} % Cancels to zero
\usepackage{empheq}
\usepackage{breqn}

\newcommand{\expect}[1]{\mathbb{E}\left[#1\right]}
\newcommand{\KL}[2]{D_{\mathrm{KL}}\!\left(#1 \,\|\, #2\right)}

% Visualization
\usepackage{graphicx} % for including images
\graphicspath{ {} } % Path to graphics folder
\usepackage{tikz}



%%% Formating
\usepackage{hyperref} % Hyperlinks
\hypersetup{
    colorlinks=true,
    linkcolor=blue,
    filecolor=magenta,      
    urlcolor=cyan,
    pdftitle={Overleaf Example},
    pdfpagemode=FullScreen,
    }
\urlstyle{same}

\usepackage{mdframed} % Framed Enviroments
\newmdenv[ 
  topline=false,
  bottomline=false,
  skipabove=\topsep,
  skipbelow=\topsep
]{sidework} %% Side-work

\usepackage{lipsum} % Lorem Ipsum example text

%%%%% ------------------ %%%%%
%%% Title
\title{DUEL}
\author{Clark Miyamoto (cm6627@nyu.edu)}
\date{\today}
\begin{document}

\maketitle

\section{Introduction}
Consider an autoregressive model which acts recursively on the tokens; e.g. it generates left to right. Now consider continuous diffusion, which defines a forward \& reverse processs.


\section{Discrete Diffusion}
Now lets attempt to create a diffusion model over a discrete state space. You start with a fully masked sequence. Then you select a position to unmask, and then choose what the token should be in that position.

\paragraph{AO-ARM Notation} In Gilad's paper, they reformulate the MDM as as a AO-ARM. Vanilla MDM's use a sequence and latents, but now you make the unmasking order $\sigma$ explicit. For example $\sigma = (\{1\}, ..., \{4\})$ does 4 steps sequential, you can also run this in parallel $(\{1,2\}, \{3,4\})$. I'll denote the partially unmasked sequence as $x_{<\sigma_t}$. 

The probability of the joint is given by
\begin{align}
	p(x,\sigma) = p(x_0, \sigma_0) \prod_{t=1}^T p(\sigma_t, \{x^{(\ell)}\}_\ell | \sigma_{1:t-1}, \{x^{(j)}\}_j)
\end{align}
where $x^{(\ell)}$ is the $\ell$'th token, but becasue you can run this in parallel, we have the $\{x^{(\ell)}\}$. The history $(\sigma_{1:t-1}, \{x^{(j)}\}_{j \in \sigma_{1:t-1}})$ uniquely determines $x_{< \sigma_t}$. Sooo
\begin{align}
	p(\sigma_t \{x^{(\ell)}\}_{\ell \in \sigma_t}
 | \sigma_{1:t-1}, \{x^{(j)}\}_{j \in \sigma_{1:t-1}}) = p(\sigma_t , \{x^{(\ell)}\}_{\ell \in \sigma_t}  | x_{< \sigma_t})
 \end{align}
 
 
 So after you go over all possible partition $\sigma$, you can marignalize to get the joint over the unmasked tokens
 \begin{align}
 	p(x) & = \sum_\sigma p(x, \sigma)\\
 	& = \sum_\sigma \prod_{t=1}^T  \pi(\sigma_t | x_{< \sigma_t}) \prod_{\ell \in \sigma_t} p(x^{(\ell)} | \ell, x_{< \sigma_t})
 \end{align}
where $\sum_\sigma$ is a partition function sum (sum over all possible partition, meansing it's combinatorial). So in previous literature, $\pi(\sigma_t | x_{< \sigma_t}) = 1/ |M|$ implicitly assume that you have uniform prior over the unmasking policy. 

To construct the training object, use ELBO (which is importance sampling + Jensen's)...
\begin{align}
	\log p(x) \geq \mathbb E_{\sigma \sim q} \left[ \log \frac{p(x,\sigma)}{q(\sigma)} \right]
\end{align}
Now plug in the joint from the previous equation, and then choose the auxiliary distribution $q$ to kill the $\pi(\cdot | x_{< \sigma_t})$. And you get
\begin{align}
	\mathcal L(\theta) = \mathbb E_{q}\left[ \sum_{t=1}^T - \log p_\theta(x^{(\sigma_t)} | x_{< \sigma_t}) \right]
\end{align}
This is the same notation $\{x^{(\ell)}\}_\ell$. So to train, I bet you use a image classifier esq neural network... So the output is the $p(x^{(\sigma_t)} | x_{< \sigma_t})$, and the input is the $x_{< \sigma_t}$. 

\section{DUEL Sampler}
So in training the unmasking is random, but at inference time it's deterministic unmasking. Gilad attempt to unite these methods by introducing a determinnistic unmasking $F: \text{tokens} \to \text{next unmasked indicies}$. He then say you can unite a bunch of unmasking rules in this way.

\begin{definition}
	[DUEL Sampler] is a pair $(x_\theta,F)$ s.t.
	\begin{itemize}
		\item Denoising network $x_\theta$ outputs a token of probability $P = \text{softmax}(x_\theta(z))$, where $z$ is the partially masked sequence.
		\item Unmasking rule $F$: deterministic map selecting positions to revel
	\end{itemize}
\end{definition}
So the algorithm goes like, while the token are not completely unmasked. Sample $P \leftarrow \text{softmax}(x)$ and then choose what to unmask $F(z)$. Now perform the unmasking by sampling $x_\ell \sim \text{Cat}(P)$ and then update the sequence of newly-partially-unmasked tokens.

\begin{theorem}
	Under deterministic policy $\pi^F$, exactly one ordered partition $\sigma^*$ has nonzero probability. The log likelihood is
	\begin{align}
		\log p^\pi = \sum_t \sum_\ell \log p(x^{(\ell)} | \ell, x_{< \sigma_t^*})
	\end{align}
	where $\sigma_t^* = F(x_{<\sigma_t^*})$. This is exact (and deterministic), which gains benefit over over ELBO.
\end{theorem}
The proof comes from the $F$ being deterministic, and hence it can be probabilistically written as a dirac delta.
















\end{document}